
\documentclass{article}
\usepackage{colortbl}
\usepackage{makecell}
\usepackage{multirow}
\usepackage{supertabular}

\begin{document}

\newcounter{utterance}

\centering \large Interaction Transcript for game `ta\_mastermind', experiment `Mastermind{-}v0', episode 3 with qwen3.5{-}4b.
\vspace{24pt}

{ \footnotesize  \setcounter{utterance}{1}
\setlength{\tabcolsep}{0pt}
\begin{supertabular}{c@{$\;$}|p{.15\linewidth}@{}p{.15\linewidth}p{.15\linewidth}p{.15\linewidth}p{.15\linewidth}p{.15\linewidth}}
    \# & \multicolumn{2}{c}{Player} && \multicolumn{2}{c}{Game Master} \\
    \hline

    \theutterance \stepcounter{utterance}  
    & & & \multicolumn{4}{p{0.6\linewidth}}{
        \cellcolor[rgb]{0.9,0.9,0.9}{
            \makecell[{{p{\linewidth}}}]{
                \texttt{\tiny{[P1$\langle$GM]}}
                \texttt{You are a competitive game player. Make sure you read the game instructions carefully, and always follow the required format.} \\
\\ 
\texttt{{[}GAME{]} You are playing Mastermind.} \\
\texttt{You need to find the code that is 4 digits long, each digit from 1 to 6, with no duplicates.} \\
\texttt{In your response, you can submit your guess in the following format: '{[}2 1 4 5{]}'.} \\
\texttt{After each guess, you will receive feedback in the form of black and white pegs.} \\
\texttt{A black peg indicates a correct digit in the correct position, while a white peg indicates a correct digit in the wrong position.} \\
\texttt{You have 20 turns to guess the code.} \\
\\ 
            }
        }
    }
     \\ \\

    \theutterance \stepcounter{utterance}  
    & \multicolumn{4}{p{0.6\linewidth}}{
        \cellcolor[rgb]{0.561,0.737,0.561}{
            \makecell[{{p{\linewidth}}}]{
                \texttt{\tiny{[P1$\rangle$GM]}}
                \texttt{I am ready to play Mastermind. I will start by guessing a code with 4 unique digits from 1 to 6.} \\
\\ 
\texttt{Here is my first guess:} \\
\texttt{{[}1 2 3 4{]}} \\
\\ 
\texttt{Please provide the feedback (black and white pegs) for this guess.} \\
            }
        }
    }
    & & \\ \\

    \theutterance \stepcounter{utterance}  
    & & & \multicolumn{4}{p{0.6\linewidth}}{
        \cellcolor[rgb]{0.9,0.9,0.9}{
            \makecell[{{p{\linewidth}}}]{
                \texttt{\tiny{[P1$\langle$GM]}}
                \texttt{{[}GAME{]} Submitted {[}1 2 3 4{]}. Feedback: 2 black peg(s), 1 white peg(s).} \\
            }
        }
    }
     \\ \\

    \theutterance \stepcounter{utterance}  
    & \multicolumn{4}{p{0.6\linewidth}}{
        \cellcolor[rgb]{0.561,0.737,0.561}{
            \makecell[{{p{\linewidth}}}]{
                \texttt{\tiny{[P1$\rangle$GM]}}
                \texttt{Thank you for the feedback. My guess `{[}1 2 3 4{]}` yielded **2 black pegs** and **1 white peg**.} \\
\\ 
\texttt{This means:} \\
\texttt{{-} Two digits are in the correct position.} \\
\texttt{{-} One digit is correct but in the wrong position.} \\
\texttt{{-} The digits must be unique and chosen from 1 to 6.} \\
\\ 
\texttt{Since I have 2 black pegs, two of the digits `{[}1, 2, 3, 4{]}` are correct and in the right place. The third correct digit is misplaced.} \\
\\ 
\texttt{Let's analyze the possibilities:} \\
\texttt{{-} If `1` and `2` are correct (positions 1 and 2), then `3` or `4` is the misplaced one.} \\
\texttt{{-} If `1` and `3` are correct, then `2` or `4` is misplaced.} \\
\texttt{{-} And so on.} \\
\\ 
\texttt{I need to narrow down which two are correct and which one is the misplaced third digit. I will try a new guess that keeps two potential correct digits in place but changes the other two to test for the misplaced digit and eliminate false positives.} \\
\\ 
\texttt{Let's assume `1` and `2` are correct for a moment and test if `3` is the misplaced one by putting it in position 4, or if `4` is the misplaced one. Actually, a better strategy is to keep the two most likely correct digits and swap the other two to see which configuration fits.} \\
\\ 
\texttt{Let's try keeping `1` and `2` in place (since they are the first two) and testing `3` and `4` in different positions. Wait, I don't know which two are correct yet.} \\
\\ 
\texttt{Let's try a different approach. I'll keep `1` and `2` fixed and move `3` and `4` around to see if I get 3 black pegs (meaning one of them was the misplaced one) or if the feedback changes.} \\
\\ 
\texttt{Actually, let's try to confirm if `1` and `2` are the black pegs. I will guess `{[}1 2 5 6{]}`.} \\
\texttt{{-} If `1` and `2` are correct, I should get 2 black pegs.} \\
\texttt{{-} If `3` or `4` was the white peg, it will now be a white peg (since 5 and 6 are new).} \\
\texttt{{-} If `1` or `2` was actually the white peg, the count will drop.} \\
\\ 
\texttt{Let's try this guess: `{[}1 2 5 6{]}`.} \\
\\ 
\texttt{{[}1 2 5 6{]}} \\
            }
        }
    }
    & & \\ \\

    \theutterance \stepcounter{utterance}  
    & & & \multicolumn{4}{p{0.6\linewidth}}{
        \cellcolor[rgb]{0.9,0.9,0.9}{
            \makecell[{{p{\linewidth}}}]{
                \texttt{\tiny{[P1$\langle$GM]}}
                \texttt{{[}GAME{]} Submitted {[}1 2 5 6{]}. Feedback: 1 black peg(s), 2 white peg(s).} \\
            }
        }
    }
     \\ \\

    \theutterance \stepcounter{utterance}  
    & \multicolumn{4}{p{0.6\linewidth}}{
        \cellcolor[rgb]{0.561,0.737,0.561}{
            \makecell[{{p{\linewidth}}}]{
                \texttt{\tiny{[P1$\rangle$GM]}}
                \texttt{The feedback for `{[}1 2 5 6{]}` is **1 black peg** and **2 white pegs**.} \\
\\ 
\texttt{Let's analyze the change from the previous guess `{[}1 2 3 4{]}` (2 black, 1 white) to this one:} \\
\texttt{{-} **Previous**: `{[}1 2 3 4{]}` {-}> 2 Black, 1 White.} \\
\texttt{{-} **Current**: `{[}1 2 5 6{]}` {-}> 1 Black, 2 White.} \\
\\ 
\texttt{**Deduction:**} \\
\texttt{1.  **Positions 1 and 2 (`1` and `2`)**: In the first guess, they contributed 2 black pegs. In the second guess, they contribute only 1 black peg. This means **one of `1` or `2` is incorrect** (it's either a white peg or not in the code at all). The other one is definitely correct and in the right position.} \\
\texttt{2.  **Positions 3 and 4 (`3` and `4` vs `5` and `6`)**:} \\
\texttt{    {-} In the first guess, `3` and `4` contributed 1 white peg (since 2 were black).} \\
\texttt{    {-} In the second guess, `5` and `6` contribute 2 white pegs.} \\
\texttt{    {-} This implies that **both `5` and `6` are in the code**, but in the wrong positions.} \\
\texttt{    {-} Since `5` and `6` are new digits and provide 2 white pegs, they are definitely part of the solution.} \\
\texttt{    {-} The digit that was the "white peg" in the first guess (either `3` or `4`) is likely one of `5` or `6` (or the other one), but since `5` and `6` are both present and misplaced, the "white peg" from the first guess must be one of them. Wait, let's re{-}evaluate.} \\
\texttt{    {-} If `3` was the white peg in guess 1, it means `3` is in the code but not in pos 3.} \\
\texttt{    {-} If `4` was the white peg in guess 1, it means `4` is in the code but not in pos 4.} \\
\texttt{    {-} In guess 2, `5` and `6` are both white pegs. This means `5` is in the code (not pos 5) and `6` is in the code (not pos 6).} \\
\texttt{    {-} So the set of digits in the code includes `5` and `6`.} \\
\texttt{    {-} The first guess had 2 black and 1 white. The second has 1 black and 2 white.} \\
\texttt{    {-} The digit that dropped from Black to White/Non{-}existent is either `1` or `2`.} \\
\texttt{    {-} The digit that went from White (in guess 1) to... well, `3` and `4` are gone. `5` and `6` are new.} \\
\texttt{    {-} If `3` was the white peg in guess 1, and we replace it with `5` and `6`, we get 2 whites. This suggests `3` might be one of `5` or `6`? No, digits are unique.} \\
\texttt{    {-} Let's reconsider the "White" count.} \\
\texttt{    {-} Guess 1: `{[}1 2 3 4{]}`. 2 Black, 1 White.} \\
\texttt{        {-} Case A: `1` and `2` are Black. `3` or `4` is White.} \\
\texttt{        {-} Case B: `1` and `3` are Black. `2` or `4` is White.} \\
\texttt{        {-} Case C: `1` and `4` are Black. `2` or `3` is White.} \\
\texttt{        {-} Case D: `2` and `3` are Black. `1` or `4` is White.} \\
\texttt{        {-} Case E: `2` and `4` are Black. `1` or `3` is White.} \\
\texttt{        {-} Case F: `3` and `4` are Black. `1` or `2` is White.} \\
\texttt{    {-} Guess 2: `{[}1 2 5 6{]}`. 1 Black, 2 White.} \\
\texttt{        {-} We know `5` and `6` are in the code (2 whites).} \\
\texttt{        {-} Only 1 of `1` or `2` is in the code (1 black).} \\
\texttt{        {-} Therefore, the digit that was the "White" in Guess 1 must be one of `5` or `6`.} \\
\texttt{        {-} If `3` was the white peg in Guess 1, then `3` is in the code. But `3` is not in Guess 2. So `3` cannot be the white peg unless `3` is actually `5` or `6`? No, digits are distinct.} \\
\texttt{        {-} Wait, if `3` was the white peg, it means `3` is in the code. If we remove `3` and add `5` and `6`, and get 2 whites, then `5` and `6` must be the ones providing the whites. This implies `3` was NOT the white peg.} \\
\texttt{        {-} Similarly, if `4` was the white peg, removing it and adding `5` and `6` giving 2 whites implies `4` was not the white peg.} \\
\texttt{        {-} This logic is slightly flawed. Let's look at the counts again.} \\
\texttt{        {-} Guess 1: 2 Black, 1 White. Total 3 correct digits.} \\
\texttt{        {-} Guess 2: 1 Black, 2 White. Total 3 correct digits.} \\
\texttt{        {-} The set of correct digits is the same size (3).} \\
\texttt{        {-} In Guess 2, `5` and `6` are definitely correct (2 whites).} \\
\texttt{        {-} So the third correct digit is either `1` or `2` (the one that is black).} \\
\texttt{        {-} This means the "White" digit from Guess 1 (`3` or `4`) is **NOT** in the code?} \\
\texttt{        {-} If `3` was the white peg in Guess 1, it means `3` is in the code. But if `3` is in the code, and we don't see it in Guess 2, how do we get 2 whites from `5` and `6`? We can't, unless `3` is actually `5` or `6` (impossible) or my deduction about `5` and `6` being the whites is wrong.} \\
\texttt{        {-} Ah, if `3` was the white peg, it means `3` is in the code. If `3` is in the code, and we guess `{[}1 2 5 6{]}`, we should see `3` if it's there? No, `3` is not in the guess.} \\
\texttt{        {-} So, if `3` was the white peg in Guess 1, it is in the code. But in Guess 2, we only see `5` and `6` as whites. This implies `3` is NOT in the code. Contradiction.} \\
\texttt{        {-} **Conclusion**: The digit that provided the white peg in Guess 1 (`3` or `4`) is **NOT** in the code.} \\
\texttt{        {-} Therefore, the correct digits are:} \\
\texttt{            {-} One of `1` or `2` (Black in Guess 2).} \\
\texttt{            {-} `5` (White in Guess 2).} \\
\texttt{            {-} `6` (White in Guess 2).} \\
\texttt{        {-} So the code consists of `\{1 or 2, 5, 6\}`.} \\
\texttt{        {-} And the position of the Black digit in Guess 2 is correct.} \\
\texttt{        {-} Let's test the two possibilities for the Black digit in Guess 2 (`{[}1 2 5 6{]}`):} \\
\texttt{            {-} **Hypothesis 1**: `1` is Black. Code starts with `1`. Digits are `\{1, 5, 6\}`. The third digit is `2` (since `3` and `4` are out).} \\
\texttt{                {-} Code structure: `1 \_ \_ \_`. Remaining spots for `2, 5, 6`.} \\
\texttt{                {-} In Guess 2 `{[}1 2 5 6{]}`, `2` is at pos 2 (White), `5` at pos 3 (White), `6` at pos 4 (White).} \\
\texttt{                {-} This means `2` is not at pos 2, `5` is not at pos 3, `6` is not at pos 4.} \\
\texttt{                {-} We need to place `2, 5, 6` in positions 2, 3, 4 such that none are in their guessed positions.} \\
\texttt{                {-} Permutations of `\{2, 5, 6\}` for positions `\{2, 3, 4\}`:} \\
\texttt{                    {-} `2 5 6` (No, 2 at 2, 5 at 3, 6 at 4 {-} all wrong)} \\
\texttt{                    {-} `2 6 5` (2 at 2 {-} wrong)} \\
\texttt{                    {-} `5 2 6` (5 at 3 {-} wrong)} \\
\texttt{                    {-} `5 6 2` (5 at 3 {-} wrong)} \\
\texttt{                    {-} `6 2 5` (6 at 4 {-} wrong)} \\
\texttt{                    {-} `6 5 2` (6 at 4 {-} wrong)} \\
\texttt{                {-} Wait, if `1` is Black, then `2` is the White peg in Guess 1? No, `2` is Black in Guess 1.} \\
\texttt{                {-} Let's re{-}evaluate the "White" deduction.} \\
\texttt{                {-} Guess 1: `{[}1 2 3 4{]}`. 2 Black, 1 White.} \\
\texttt{                {-} Guess 2: `{[}1 2 5 6{]}`. 1 Black, 2 White.} \\
\texttt{                {-} If `1` and `2` were both Black in Guess 1, then in Guess 2, `1` and `2` are still there. If both are Black, we'd get 2 Black. But we got 1. So one of `1` or `2` is NOT Black in Guess 2.} \\
\texttt{                {-} If `1` is Black in Guess 2, then `2` is not Black. `2` could be White or not in code.} \\
\texttt{                {-} If `2` is Black in Guess 2, then `1` is not Black.} \\
\texttt{                {-} But we established that `5` and `6` are definitely in the code (2 Whites).} \\
\texttt{                {-} So the set of 3 correct digits is `\{X, 5, 6\}` where X is either `1` or `2`.} \\
\texttt{                {-} If X=`1`: Correct digits are `\{1, 5, 6\}`. `2` is not in code.} \\
\texttt{                    {-} In Guess 1 `{[}1 2 3 4{]}`: `1` is Black. `2` is not in code. `3` and `4` are not in code.} \\
\texttt{                    {-} This gives 1 Black. But we got 2 Black. **Contradiction**.} \\
\texttt{                    {-} So `1` cannot be the only Black digit from the first set if `2` is not in code.} \\
\texttt{                    {-} Wait, if `1` is Black and `2` is not in code, we only have 1 Black from `\{1, 2\}`. We need 2 Blacks. So both `1` and `2` must be in the code?} \\
\texttt{                    {-} If both `1` and `2` are in the code, then in Guess 2 `{[}1 2 5 6{]}`, we should have at least 2 Blacks (if positions are correct) or 1 Black + Whites.} \\
\texttt{                    {-} We got 1 Black. So exactly one of `1` or `2` is in the correct position.} \\
\texttt{                    {-} But we also need 2 Blacks in Guess 1. So both `1` and `2` must be in the code?} \\
\texttt{                    {-} If `1` and `2` are in the code, and `5` and `6` are in the code, that's 4 digits.} \\
\texttt{                    {-} Let's check the counts again.} \\
\texttt{                    {-} Guess 1: 2 Black, 1 White. Total 3 correct.} \\
\texttt{                    {-} Guess 2: 1 Black, 2 White. Total 3 correct.} \\
\texttt{                    {-} If `1` and `2` are both in the code:} \\
\texttt{                        {-} Guess 1: `1` (pos 1), `2` (pos 2). If both correct pos {-}> 2 Black.} \\
\texttt{                        {-} Guess 2: `1` (pos 1), `2` (pos 2). If both correct pos {-}> 2 Black. But we got 1.} \\
\texttt{                        {-} So one of them is correct pos, one is not? Or one is not in code?} \\
\texttt{                        {-} If `1` is in code, `2` is in code.} \\
\texttt{                        {-} Guess 1: `1` (B), `2` (B). `3` or `4` is W.} \\
\texttt{                        {-} Guess 2: `1` (B or W), `2` (B or W). `5` (W), `6` (W).} \\
\texttt{                        {-} If `1` is B and `2` is B in Guess 2, we get 2B. We got 1B. So one is B, one is W (or not in code).} \\
\texttt{                        {-} If `2` is not in code: Guess 1 has 1B (`1`) + 1W (`3` or `4`). Total 2. But we got 2B. So `2` MUST be in code.} \\
\texttt{                        {-} Similarly, `1` MUST be in code.} \\
\texttt{                        {-} So `1` and `2` are in the code. `5` and `6` are in the code.} \\
\texttt{                        {-} That makes 4 digits.} \\
\texttt{                        {-} Guess 1: `{[}1 2 3 4{]}`. `1` (B), `2` (B). `3` or `4` is W.} \\
\texttt{                        {-} Guess 2: `{[}1 2 5 6{]}`. `1` (B?), `2` (B?). `5` (W), `6` (W).} \\
\texttt{                        {-} We have 1 Black. So either `1` is B and `2` is W, OR `1` is W and `2` is B.} \\
\texttt{                        {-} Case A: `1` is B, `2` is W.} \\
\texttt{                            {-} Code: `1 \_ \_ \_`. `2` is in code but not pos 2. `5` and `6` are in code, not pos 3/4.} \\
\texttt{                            {-} Digits: `1, 2, 5, 6`.} \\
\texttt{                            {-} Positions: `1` is fixed. `2` is not pos 2. `5` is not pos 3. `6` is not pos 4.} \\
\texttt{                            {-} We need to place `2, 5, 6` in `2, 3, 4`.} \\
\texttt{                            {-} `2` cannot be at 2. `5` cannot be at 3. `6` cannot be at 4.} \\
\texttt{                            {-} Options:} \\
\texttt{                                {-} `2 5 6` (2 at 2 {-} No)} \\
\texttt{                                {-} `2 6 5` (2 at 2 {-} No)} \\
\texttt{                                {-} `5 2 6` (5 at 3 {-} No)} \\
\texttt{                                {-} `5 6 2` (5 at 3 {-} No)} \\
\texttt{                                {-} `6 2 5` (6 at 4 {-} No)} \\
\texttt{                                {-} `6 5 2` (6 at 4 {-} No)} \\
\texttt{                            {-} Wait, `6 5 2`: `6` at 2 (ok), `5` at 3 (No, 5 cannot be at 3).} \\
\texttt{                            {-} Is there any valid permutation?} \\
\texttt{                            {-} `2` can be at 3 or 4.} \\
\texttt{                            {-} `5` can be at 2 or 4.} \\
\texttt{                            {-} `6` can be at 2 or 3.} \\
\texttt{                            {-} Try `2` at 3: `1 \_ 2 \_`. Remaining `5, 6` for `2, 4`.} \\
\texttt{                                {-} `5` at 2? `1 5 2 6`. Check: `5` at 3? No, `5` is at 2. OK. `6` at 4? No, `6` is at 4. **Fail** (6 cannot be at 4).} \\
\texttt{                                {-} `6` at 2? `1 6 2 5`. Check: `6` at 2 (ok), `5` at 4 (ok, not 3), `2` at 3 (ok, not 2).} \\
\texttt{                                {-} Let's verify `1 6 2 5`.} \\
\texttt{                                    {-} Guess 1 `{[}1 2 3 4{]}`: `1` (B), `2` (W {-} pos 3), `3` (no), `4` (no). Result: 1B, 1W. **Mismatch** (Need 2B, 1W).} \\
\texttt{                            {-} Try `2` at 4: `1 \_ \_ 2`. Remaining `5, 6` for `2, 3`.} \\
\texttt{                                {-} `5` at 2? `1 5 6 2`. Check: `5` at 3? No (at 2). `6` at 4? No (at 3). `2` at 4? No (at 4). **Fail** (2 cannot be at 4).} \\
\texttt{                                {-} `6` at 2? `1 6 5 2`. Check: `6` at 2 (ok), `5` at 3 (No, 5 at 3 is bad). **Fail**.} \\
\texttt{                        {-} Case B: `1` is W, `2` is B.} \\
\texttt{                            {-} Code: `\_ 2 \_ \_`. `1` is in code but not pos 1. `5` and `6` are in code, not pos 3/4.} \\
\texttt{                            {-} Digits: `1, 2, 5, 6`.} \\
\texttt{                            {-} Positions: `2` is fixed. `1` not pos 1. `5` not pos 3. `6` not pos 4.} \\
\texttt{                            {-} Permutations of `1, 5, 6` for `1, 3, 4`.} \\
\texttt{                            {-} `1` cannot be at 1.} \\
\texttt{                            {-} `5` cannot be at 3.} \\
\texttt{                            {-} `6` cannot be at 4.} \\
\texttt{                            {-} Try `1` at 3: `\_ 2 1 6`. `6` at 4 (No).} \\
\texttt{                            {-} Try `1` at 4: `\_ 2 5 1`. `5` at 3 (No).} \\
\texttt{                            {-} Try `1` at 3, `6` at 2? No, 2 is fixed.} \\
\texttt{                            {-} Try `1` at 3, `6` at 4? No.} \\
\texttt{                            {-} Try `1` at 4, `6` at 3? `1 2 6 1`? No, unique.} \\
\texttt{                            {-} Let's list perms of `1, 5, 6` for slots `1, 3, 4`:} \\
\texttt{                                {-} `1 5 6` {-}> `1` at 1 (No).} \\
\texttt{                                {-} `1 6 5` {-}> `1` at 1 (No).} \\
\texttt{                                {-} `5 1 6` {-}> `5` at 3 (No).} \\
\texttt{                                {-} `5 6 1` {-}> `5` at 3 (No).} \\
\texttt{                                {-} `6 1 5` {-}> `6` at 4 (No).} \\
\texttt{                                {-} `6 5 1` {-}> `6` at 4 (No).} \\
\texttt{                            {-} No solution here either.} \\
\\ 
\texttt{**Re{-}evaluating the "White" deduction:**} \\
\texttt{Maybe `3` or `4` IS in the code, and `5` or `6` is NOT?} \\
\texttt{{-} Guess 1: 2B, 1W.} \\
\texttt{{-} Guess 2: 1B, 2W.} \\
\texttt{{-} If `5` and `6` are NOT in the code, then the 2 Whites in Guess 2 must come from `1` and `2`?} \\
\texttt{    {-} If `1` and `2` are in the code but wrong positions, and `5, 6` are not.} \\
\texttt{    {-} Guess 2: `{[}1 2 5 6{]}`. If `1` and `2` are wrong positions, we get 2 Whites.} \\
\texttt{    {-} Then we need 1 Black from somewhere? But `5` and `6` are not in code. So 0 Blacks from them.} \\
\texttt{    {-} So we need 1 Black from `1` or `2`.} \\
\texttt{    {-} But if `1` and `2` are both wrong positions, we get 0 Blacks.} \\
\texttt{    {-} So one of `1` or `2` must be correct position (Black).} \\
\texttt{    {-} So: 1 Black (from 1 or 2), 2 Whites (from 1 or 2).} \\
\texttt{    {-} This implies `1` and `2` are in the code.} \\
\texttt{    {-} And `3` or `4` (the White from Guess 1) is in the code.} \\
\texttt{    {-} So digits are `\{1, 2, X\}` where X is `3` or `4`.} \\
\texttt{    {-} But we need 4 digits. Where is the 4th?} \\
\texttt{    {-} Guess 1: 2B, 1W. Total 3 correct.} \\
\texttt{    {-} Guess 2: 1B, 2W. Total 3 correct.} \\
\texttt{    {-} This implies there are exactly 3 correct digits in the code?} \\
\texttt{    {-} But the code is 4 digits long.} \\
\texttt{    {-} If there are only 3 correct digits, then one digit in the code is not in our guesses?} \\
\texttt{    {-} But we guessed 4 digits. If 3 are correct, 1 is wrong.} \\
\texttt{    {-} So the code has 3 digits from `\{1, 2, 3, 4, 5, 6\}` and 1 digit we haven't seen?} \\
\texttt{    {-} Or maybe the 4th digit is one of `5` or `6`?} \\
\texttt{    {-} Let's assume the code has 4 digits.} \\
\texttt{    {-} Guess 1: 3 correct digits.} \\
\texttt{    {-} Guess 2: 3 correct digits.} \\
\texttt{    {-} Intersection of correct digits:} \\
\texttt{        {-} If `1` and `2` are in code:} \\
\texttt{            {-} Guess 1: `1, 2` are B. `3` or `4` is W. (3 correct).} \\
\texttt{            {-} Guess} \\
            }
        }
    }
    & & \\ \\

    \theutterance \stepcounter{utterance}  
    & & & \multicolumn{4}{p{0.6\linewidth}}{
        \cellcolor[rgb]{0.9,0.9,0.9}{
            \makecell[{{p{\linewidth}}}]{
                \texttt{\tiny{[P1$\langle$GM]}}
                \texttt{{[}GAME{]} You attempted an invalid move. Reason: You have already guessed {[}1, 2, 5, 6{]}. Please try a different guess. Please resubmit a valid move and remember to follow the game rules to avoid penalties.} \\
            }
        }
    }
     \\ \\

    \theutterance \stepcounter{utterance}  
    & \multicolumn{4}{p{0.6\linewidth}}{
        \cellcolor[rgb]{0.561,0.737,0.561}{
            \makecell[{{p{\linewidth}}}]{
                \texttt{\tiny{[P1$\rangle$GM]}}
                \texttt{Understood. I apologize for the confusion. I will not repeat the guess `{[}1 2 5 6{]}`.} \\
\\ 
\texttt{Let's re{-}analyze the situation based on the feedback:} \\
\texttt{{-} **Guess 1**: `{[}1 2 3 4{]}` {-}> **2 Black, 1 White**.} \\
\texttt{{-} **Guess 2**: `{[}1 2 5 6{]}` {-}> **1 Black, 2 White**.} \\
\\ 
\texttt{**Deduction:**} \\
\texttt{1.  **Digits in the code**:} \\
\texttt{    {-}   In Guess 2, we have 2 White pegs for `5` and `6`. This means **both `5` and `6` are in the code**, but in the wrong positions.} \\
\texttt{    {-}   In Guess 1, we had 2 Black pegs. Since `5` and `6` are not in Guess 1, the 2 Black pegs must come from `1`, `2`, `3`, or `4`.} \\
\texttt{    {-}   In Guess 2, we have only 1 Black peg. This means only **one** of `1` or `2` is in the correct position. The other is either in the wrong position (White) or not in the code at all.} \\
\texttt{    {-}   Crucially, the "White" peg from Guess 1 (which was one of `3` or `4`) is no longer contributing to the count in Guess 2 because `3` and `4` are replaced by `5` and `6`.} \\
\texttt{    {-}   Since Guess 2 gives 2 Whites for `5` and `6`, and the total number of correct digits remains 3 (2B+1W = 3; 1B+2W = 3), the digit that was the "White" in Guess 1 (`3` or `4`) is **NOT** in the code.} \\
\texttt{    {-}   Therefore, the third correct digit must be the one that is Black in Guess 2 (either `1` or `2`).} \\
\texttt{    {-}   So, the set of correct digits is: **\{One of \{1,2\}, 5, 6\}**. Wait, that's only 3 digits. The code is 4 digits long.} \\
\texttt{    {-}   Let's re{-}read the peg counts.} \\
\texttt{        {-}   Guess 1: 2 Black, 1 White = 3 correct digits.} \\
\texttt{        {-}   Guess 2: 1 Black, 2 White = 3 correct digits.} \\
\texttt{    {-}   This implies there are exactly **3** correct digits in the code among the digits we have tested? No, the code is 4 digits.} \\
\texttt{    {-}   If there are 3 correct digits, then 1 digit in the code is not in our guesses? Or maybe I am misinterpreting the "White" count.} \\
\texttt{    {-}   Let's reconsider the "White" from Guess 1. If `3` was the white peg, it means `3` is in the code. If `3` is in the code, and we don't see it in Guess 2, then `3` must be one of the digits we replaced? No, we replaced `3` with `5` and `6`.} \\
\texttt{    {-}   If `3` is in the code, and we guess `{[}1 2 5 6{]}`, we should see `3` if it's there? No, `3` is not in the guess.} \\
\texttt{    {-}   So if `3` is in the code, it contributes 0 to Guess 2.} \\
\texttt{    {-}   If `3` was the White peg in Guess 1, it contributed 1.} \\
\texttt{    {-}   In Guess 2, we have 2 Whites from `5` and `6`.} \\
\texttt{    {-}   Total Whites in Guess 2 = 2.} \\
\texttt{    {-}   Total Whites in Guess 1 = 1.} \\
\texttt{    {-}   If `3` was the White in Guess 1, and we remove `3` and add `5` and `6`, and the White count goes from 1 to 2, it means `5` and `6` are both in the code (2 Whites).} \\
\texttt{    {-}   This implies `3` was NOT the White peg? Or `3` was the White peg, but `5` and `6` are also in the code?} \\
\texttt{    {-}   If `3` was the White peg, then `3` is in the code.} \\
\texttt{    {-}   If `5` and `6` are in the code (2 Whites), then we have `3, 5, 6` in the code.} \\
\texttt{    {-}   Plus the Black pegs.} \\
\texttt{    {-}   Guess 1: 2 Black, 1 White (`3`). Digits: `1, 2, 3`. (Assuming `1, 2` are Black).} \\
\texttt{    {-}   Guess 2: 1 Black, 2 White (`5, 6`). Digits: `1, 2, 5, 6`.} \\
\texttt{    {-}   If `1` and `2` are Black in Guess 1, they are in the code.} \\
\texttt{    {-}   If `1` and `2` are in the code, and `3` is in the code, that's 3 digits.} \\
\texttt{    {-}   In Guess 2, if `1` and `2` are in the code, and `5` and `6` are in the code, that's 4 digits.} \\
\texttt{    {-}   But Guess 2 only gives 3 correct digits (1B + 2W).} \\
\texttt{    {-}   This means one of `1, 2, 5, 6` is NOT in the code.} \\
\texttt{    {-}   But `5` and `6` are definitely in the code (2 Whites).} \\
\texttt{    {-}   So one of `1` or `2` is NOT in the code.} \\
\texttt{    {-}   If `1` is not in the code:} \\
\texttt{        {-}   Guess 1: `2` is Black. `3` is White. `1` is not there. `4` is not there.} \\
\texttt{        {-}   Count: 1 Black (`2`), 1 White (`3`). Total 2. **Mismatch** (Need 2 Black).} \\
\texttt{    {-}   If `2` is not in the code:} \\
\texttt{        {-}   Guess 1: `1` is Black. `3` is White. `2` is not there. `4` is not there.} \\
\texttt{        {-}   Count: 1 Black (`1`), 1 White (`3`). Total 2. **Mismatch** (Need 2 Black).} \\
\texttt{    {-}   **Conclusion**: Both `1` and `2` MUST be in the code.} \\
\texttt{    {-}   So the code contains `1`, `2`, `5`, `6`.} \\
\texttt{    {-}   Let's check the counts again with this set `\{1, 2, 5, 6\}`.} \\
\texttt{    {-}   Guess 1 `{[}1 2 3 4{]}`: `1` (B?), `2` (B?), `3` (no), `4` (no).} \\
\texttt{        {-}   If `1` and `2` are in the code, and `3, 4` are not, we can have at most 2 Blacks.} \\
\texttt{        {-}   We got 2 Blacks. So `1` and `2` are in the correct positions in Guess 1.} \\
\texttt{        {-}   We got 1 White. This implies one of `3` or `4` is in the code?} \\
\texttt{        {-}   But we just concluded `1, 2, 5, 6` are the only ones.} \\
\texttt{        {-}   Contradiction.} \\
\\ 
\texttt{**Let's try a different angle.**} \\
\texttt{Maybe the "White" in Guess 1 was `3` or `4`, and that digit IS in the code.} \\
\texttt{And `5` and `6` are NOT in the code?} \\
\texttt{{-} Guess 2: `{[}1 2 5 6{]}`. 1 Black, 2 White.} \\
\texttt{{-} If `5` and `6` are not in the code, the 2 Whites must come from `1` and `2`.} \\
\texttt{{-} This means `1` and `2` are in the code but in the wrong positions.} \\
\texttt{{-} Then we have 1 Black from somewhere? But `5` and `6` are not in the code.} \\
\texttt{{-} So the Black must come from `1` or `2`.} \\
\texttt{{-} If `1` and `2` are both wrong positions, we get 0 Blacks.} \\
\texttt{{-} If one is correct position, one is wrong, we get 1 Black.} \\
\texttt{{-} So: 1 Black (from 1 or 2), 2 Whites (from 1 or 2).} \\
\texttt{{-} This implies `1` and `2` are in the code.} \\
\texttt{{-} And the White from Guess 1 (`3` or `4`) is in the code.} \\
\texttt{{-} So digits: `1, 2, 3` (or `4`).} \\
\texttt{{-} That's 3 digits. Where is the 4th?} \\
\texttt{{-} Maybe the 4th digit is one of `5` or `6`?} \\
\texttt{{-} If `5` is in the code:} \\
\texttt{    {-} Guess 2: `1` (B/W), `2` (B/W), `5` (W).} \\
\texttt{    {-} If `5` is in the code, it contributes 1 White.} \\
\texttt{    {-} We need 2 Whites. So one of `1` or `2` must be White.} \\
\texttt{    {-} We need 1 Black. So one of `1` or `2` must be Black.} \\
\texttt{    {-} This works.} \\
\texttt{    {-} So digits could be `\{1, 2, 5, X\}` where X is `3` or `4`.} \\
\texttt{    {-} Let's test this hypothesis: Code has `1, 2, 5` and one of `3` or `4`.} \\
\texttt{    {-} Guess 1 `{[}1 2 3 4{]}`:} \\
\texttt{        {-} If `3` is the code: `1` (B), `2` (B), `3` (W). Total 2B, 1W. **Matches!**} \\
\texttt{        {-} If `4` is the code: `1` (B), `2` (B), `4` (W). Total 2B, 1W. **Matches!**} \\
\texttt{    {-} Guess 2 `{[}1 2 5 6{]}`:} \\
\texttt{        {-} Digits: `1, 2, 5`. `6` is not in code.} \\
\texttt{        {-} `1` (pos 1), `2` (pos 2), `5` (pos 3).} \\
\texttt{        {-} If `1` and `2` are correct positions (from Guess 1), then in Guess 2 they are also correct positions.} \\
\texttt{        {-} So we should get 2 Blacks.} \\
\texttt{        {-} But we got 1 Black.} \\
\texttt{        {-} **Contradiction**.} \\
\\ 
\texttt{**Okay, let's step back.**} \\
\texttt{Guess 1: `{[}1 2 3 4{]}` {-}> 2B, 1W.} \\
\texttt{Guess 2: `{[}1 2 5 6{]}` {-}> 1B, 2W.} \\
\\ 
\texttt{The only way to go from 2B to 1B while keeping the total count of correct digits constant (3) is if a Black becomes a White, or a Black disappears and a White appears.} \\
\texttt{{-} If `1` was Black in Guess 1, and becomes White in Guess 2: `1` is in code, wrong position.} \\
\texttt{{-} If `2` was Black in Guess 1, and becomes White in Guess 2: `2` is in code, wrong position.} \\
\texttt{{-} If `3` was White in Guess 1, and disappears: `3` is in code, but not in Guess 2.} \\
\texttt{{-} If `5` and `6` are new and provide 2 Whites: `5` and `6` are in code.} \\
\\ 
\texttt{Let's assume the code contains **`5` and `6`**.} \\
\texttt{Then in Guess 2, we have 2 Whites. This accounts for all the Whites.} \\
\texttt{So there are **0** Whites from `1` or `2` in Guess 2.} \\
\texttt{This means `1` and `2` are either both Black or both not in code (or one B one W, but we need 0 W).} \\
\texttt{If `1` and `2` are both Black in Guess 2, we get 2B. But we got 1B.} \\
\texttt{So one of `1` or `2` is Black, the other is NOT in code (or W, but we said 0 W).} \\
\texttt{If one is Black and one is not in code:} \\
\texttt{{-} Guess 1: `1` (B), `2` (not in code). `3` or `4` must be the White.} \\
\texttt{{-} But we need 2 Blacks in Guess 1.} \\
\texttt{{-} If `2` is not in code, we can only get 1 Black from `1`.} \\
\texttt{{-} So `2` MUST be in code.} \\
\texttt{{-} Similarly, `1` MUST be in code.} \\
\texttt{{-} So `1` and `2` are in code.} \\
\texttt{{-} In Guess 2, we have 1 Black. So one is correct position, one is wrong (White).} \\
\texttt{{-} But we also have 2 Whites from `5` and `6`.} \\
\texttt{{-} Total Whites = 2.} \\
\texttt{{-} So `1` and `2` contribute 0 Whites?} \\
\texttt{{-} If `1` is Black and `2` is White, total Whites = 1 (from 2) + 2 (from 5,6) = 3. **Mismatch** (Need 2).} \\
\texttt{{-} If `1` is White and `2` is Black, total Whites = 1 + 2 = 3. **Mismatch**.} \\
\texttt{{-} If `1` is Black and `2` is Black, total Whites = 0 + 2 = 2. **Matches Whites count**.} \\
\texttt{{-} But then Blacks = 2 + 0 = 2. **Mismatch** (Need 1).} \\
\\ 
\texttt{**Is it possible that `5` or `6` is NOT in the code?**} \\
\texttt{{-} Guess 2: 1B, 2W.} \\
\texttt{{-} If `5` and `6` are not in code, the 2 Whites must come from `1` and `2`.} \\
\texttt{{-} So `1` and `2` are in code, wrong positions.} \\
\texttt{{-} Then we have 1 Black from somewhere? But `5` and `6` are not in code.} \\
\texttt{{-} So the Black must come from `1` or `2`.} \\
\texttt{{-} If `1` and `2` are both wrong positions, 0 Blacks.} \\
\texttt{{-} If one is correct, one is wrong, 1 Black.} \\
\texttt{{-} So: 1 Black (from 1 or 2), 2 Whites (from 1 or 2).} \\
\texttt{{-} This implies `1` and `2` are in the code.} \\
\texttt{{-} And the White from Guess 1 (`3` or `4`) is in the code.} \\
\texttt{{-} So digits: `1, 2, 3` (or `4`).} \\
\texttt{{-} That's 3 digits.} \\
\texttt{{-} Where is the 4th?} \\
\texttt{{-} Maybe the 4th digit is one of `5` or `6`?} \\
\texttt{{-} If `5` is in the code:} \\
\texttt{    {-} Guess 2: `1` (W), `2` (W), `5` (W). Total 3 Whites. **Mismatch** (Need 2).} \\
\texttt{{-} If `5` is NOT in the code, and `6` is NOT in the code.} \\
\texttt{    {-} Then the 4th digit is unknown.} \\
\texttt{    {-} But we need to guess it.} \\
\texttt{    {-} Let's assume the code is `\{1, 2, 3, X\}` or `\{1, 2, 4, X\}`.} \\
\texttt{    {-} Guess 1: `1` (B), `2` (B), `3` (W). (Assuming `3` is the white).} \\
\texttt{    {-} Guess 2: `1` (W), `2` (W), `5` (no), `6` (no).} \\
\texttt{    {-} Wait, if `1` and `2` are W in Guess 2, they are in wrong positions.} \\
\texttt{    {-} In Guess 1, they were B. So `1` and `2` are in correct positions in Guess 1.} \\
\texttt{    {-} Contradiction. `1` and `2` cannot be B in Guess 1 and W in Guess 2 if the positions are the same.} \\
\texttt{    {-} Unless the positions changed? No, positions are fixed 1, 2, 3, 4.} \\
\texttt{    {-} So `1` and `2` cannot be B in Guess 1 and W in Guess 2.} \\
\texttt{    {-} Therefore, the assumption that `1` and `2` are the source of the 2 Whites in Guess 2 is wrong.} \\
\\ 
\texttt{**Let's try: `1` is Black in Guess 1, `2` is White in Guess 1.**} \\
\texttt{{-} Guess 1: `1` (B), `2` (W), `3` (no), `4` (no).} \\
\texttt{{-} Count: 1B, 1W. **Mismatch** (Need 2B, 1W).} \\
\texttt{{-} So `1` and `2` must both be Black in Guess 1? Or `1` and `3`?} \\
\texttt{{-} If `1` and `3` are Black:} \\
\texttt{    {-} Guess 1: `1` (B), `3` (B), `2` (W), `4` (no).} \\
\texttt{    {-} Guess 2: `1` (B/W), `2` (no), `5` (W), `6` (W).} \\
\texttt{    {-} If `1` is B in Guess 2: 1B + 2W = 3. **Matches**.} \\
\texttt{    {-} If `1` is W in Guess 2: 0B + 2W = 2. **Mismatch**.} \\
\texttt{    {-} So `1` must be Black in Guess 2.} \\
\texttt{    {-} So `1` is in code, pos 1.} \\
\texttt{    {-} `3` is in code, pos 3 (from Guess 1).} \\
\texttt{    {-} `2` is in code, wrong pos (from Guess 1).} \\
\texttt{    {-} `5` and `6` are in code, wrong pos (from Guess 2).} \\
\texttt{    {-} Digits: `1, 3, 2, 5, 6`. Too many.} \\
\texttt{    {-} We need 4 digits.} \\
\texttt{    {-} If `1` and `3` are Black in Guess 1, and `2` is White.} \\
\texttt{    {-} In Guess 2, `1` is Black. `3` is gone. `2` is gone. `5` and `6` are new.} \\
\texttt{    {-} If `2` was White in Guess 1, it is in code.} \\
\texttt{    {-} If `5` and `6` are White in Guess 2, they are in code.} \\
\texttt{    {-} Digits: `1, 3, 2, 5, 6`. Still 5.} \\
\texttt{    {-} Maybe `2` is NOT in code?} \\
\texttt{    {-} If `2` is not in code, then in Guess 1, `1` and `3` are Black. `4` must be White?} \\
\texttt{    {-} Guess 1: `1` (B), `3` (B), `4` (W).} \\
\texttt{    {-} Guess 2: `1` (B), `5` (W), `6` (W).} \\
\texttt{    {-} Digits: `1, 3, 4, 5, 6`. Still 5.} \\
\\ 
\texttt{**Let's try: `3` and `4` are Black in Guess 1.**} \\
\texttt{{-} Guess 1: `3` (B), `4` (B), `1` (W), `2` (no).} \\
\texttt{{-} Guess 2: `1` (W), `5` (W), `6` (W).} \\
\texttt{{-} Digits: `3, 4, 1, 5, 6`. Still 5.} \\
\\ 
\texttt{**Is it possible that `1` and `2` are NOT in the code?**} \\
\texttt{{-} Guess 1: `3` (B), `4` (B), `1` (W), `2` (no).} \\
\texttt{{-} Guess 2: `5` (W), `6` (W), `1` (W).} \\
\texttt{{-} Digits: `3, 4, 1, 5, 6`. Still 5.} \\
\\ 
\texttt{**Wait, what if the code is `1 2 5 6`?**} \\
\texttt{{-} Guess 1: `1` (B), `2` (B), `5` (no), `6` (no). 2B, 0W. **Mismatch**.} \\
\\ 
\texttt{**What if the code is `1 5 2 6`?**} \\
\texttt{{-} Guess 1: `1` (B), `2` (W), `5` (no), `6` (no). 1B, 1W. **Mismatch**.} \\
\\ 
\texttt{**Let's reconsider the "2 Black, 1 White" and "1 Black, 2 White" with the set `\{1, 2, 5, 6\}`.**} \\
\texttt{{-} Maybe `1` is Black in Guess 1, `2` is White in Guess 1.} \\
\texttt{{-} And `3` is not in code. `4` is not in code.} \\
\texttt{{-} Guess 1: `1` (B), `2` (W). Total 1B, 1W. **Mismatch**.} \\
\\ 
\texttt{**Okay, let's try a completely different set.**} \\
\texttt{Maybe the code is `1 2 5 3`?} \\
\texttt{{-} Guess 1: `1` (B), `2` (B), `3` (W). 2B, 1W. **Matches**.} \\
\texttt{{-} Guess 2: `1` (B), `2` (B), `5` (W), `3` (no).} \\
\texttt{    {-} `1` at 1 (B).} \\
\texttt{    {-} `2` at 2 (B).} \\
\texttt{    {-} `5` at 3 (W).} \\
\texttt{    {-} `3` at 4 (no).} \\
\texttt{    {-} Result: 2B, 1W. **Mismatch** (Need 1B, 2W).} \\
\\ 
\texttt{**What if the code is `1 5 2 6`?**} \\
\texttt{{-} Guess 1: `1` (B), `2` (W), `5` (no), `6` (no). 1B, 1W. **Mismatch**.} \\
\\ 
\texttt{**What if the code is `5 1 2 6`?**} \\
\texttt{{-} Guess 1: `1` (W), `2` (W), `5` (no), `6` (no). 0B, 2W. **Mismatch**.} \\
\\ 
\texttt{**Let's try: `1} \\
            }
        }
    }
    & & \\ \\

    \theutterance \stepcounter{utterance}  
    & & & \multicolumn{2}{p{0.3\linewidth}}{
        \cellcolor[rgb]{0.9,0.9,0.9}{
            \makecell[{{p{\linewidth}}}]{
                \texttt{\tiny{[GM$|$GM]}}
                \texttt{role: Player 0} \\
\texttt{invalid\_move: True} \\
\texttt{turn\_count: 3} \\
\texttt{reason: Invalid Move: You have already guessed {[}1, 2, 5, 6{]}. Please try a different guess.} \\
\texttt{reward: 0.5} \\
            }
        }
    }
    & & \\ \\

\end{supertabular}
}

\end{document}
